%% ma: L12-B01-0025
%% lop: 12
%% chuong: 1
%% bai: 1
%% dang: 12.1.17
%% loai: TN4
%% muc: NB
%% dap_an: "D"
%% nguon: "(NBV)-ĐỀ SỐ 15-MỖI NGÀY 1 ĐỀ THI 2025.pdf (D15, MỖI NGÀY 1 ĐỀ - Đề số 15), Phần I Câu 6"
%% kien_thuc: "Bài 1 (tính đơn điệu của hàm trùng phương)"
%% trang_thai: cho-duyet
%% ly_do: "Dạng mới 12.1.17 đề xuất (đơn điệu của hàm trùng phương y = x^4 + c, xét khẳng định đúng)."
%% ghi_chu: "Đáp án gốc D đúng."
\begin{ex}
Cho hàm số $y=x^4+5$. Khẳng định nào sau đây đúng?
\choice{Hàm số đồng biến trên $\mathbb R$}{Hàm số nghịch biến trên $\mathbb R$}{Hàm số nghịch biến trên $(-\infty;5)$}{\True Hàm số nghịch biến trên $(-\infty;0)$}
\loigiai{$y'=4x^3$, $y'<0$ khi $x<0$ và $y'>0$ khi $x>0$. Hàm số nghịch biến trên $(-\infty;0)$ và đồng biến trên $(0;+\infty)$.}
\end{ex}
